\section{Vector Calculus}
We call functions $\mathbf{r}: \mathbb{R}\to\mathbb{R}^{3}$ a \textbf{vector
function} or a \textbf{vector-valued function}.
We may write $\mathbf{r}$ as:
\[\mathbf{r}\left( t \right) =\left<f\left( t \right) ,g\left( t \right)
,h\left( t \right)  \right> =f\left( t \right) \mathbf{\hat{\imath}}+g\left( t
\right) \mathbf{\hat{\jmath}}+h\left( t \right) \hat{k}.\] 
We may interpret $\mathbf{r}\left( t \right) $ as a point $P=\left( f\left( t
\right) ,g\left( t \right) ,h\left( t \right)  \right) $ moving in space over
time $t$. We form a \textcolor{Blue}{curve} in space with its path.

\paragraph{Limits}
For vector function $\mathbf{r}\left( t \right) =f\left( t \right)
\hat{\imath}+g\left( t \right) \hat{\jmath}+h\left( t \right) \hat{k}$ and
vector $\mathbf{L}$, $\lim_{t \to t_0} \mathbf{r}\left( t \right)
=\mathbf{L}\iff\forall \epsilon>0\exists\delta>0\text{ such that }$
\[\forall t, 0<\left| t-t_0 \right| <\delta \implies\left| \mathbf{r}\left( t \right) -\mathbf{L} \right| <\epsilon.\] 
It can be shown that the limit of a vector function comprises the limits of its
coordinates:
\[\lim_{t \to t_0} \mathbf{r}\left( t \right) =\left<\lim_{t \to t_0} f\left( t \right) ,\lim_{t \to t_0} g\left( t \right) ,\lim_{t \to t_0} h\left( t \right)  \right>.\] 

The vector function is \textcolor{Bittersweet}{continuous at a point $t=t_0$}
in its domain if
\[\lim_{t \to t_0} \mathbf{r}\left( t \right) =\mathbf{r}\left( t \right)\] 
and it is a \textcolor{Bittersweet}{continuous function} $\iff$ it is
continuous at every point in its domain.

\paragraph{Differentiation}
We define the derivative of a vector function as
\[\mathbf{r}'\left( t \right) =\frac{d \mathbf{r}}{dt}=\lim_{\Delta t \to 0} \frac{\mathbf{r}\left( t+\Delta t \right) -\mathbf{r}\left( t \right) }{\Delta t}.\]
\[\mathbf{r}'\left( t \right) =\frac{df}{dt}\hat{\imath}+\frac{dg}{dt}\hat{\jmath}+\frac{dh}{dt}\hat{k}.\] 
The derivative does not exist if the limit does not exist.
We may interpret a nonzero derivative as a \textbf{tangent
vector} to the curve at the given $t$. We then have the
\textcolor{Bittersweet}{tangent line} that which is spanned by
$\mathbf{r}'\left( t \right) $.\\
If $\forall t, \mathbf{r}'\left( t \right)$ exists,
$\mathbf{r}\left( t \right) $ is \textcolor{Bittersweet}{differentiable}. If
$\mathbf{r}(t)$ is differentiable (continuous) and $\forall t,\mathbf{r}'\left(
t \right)\neq \mathbf{0}$,the curve is \textcolor{Bittersweet}{smooth}.\\
We may also interpret derivatives as the \textbf{velocity} of the point moving
in space over time. That is, $\textcolor{Bittersweet}{\mathbf{v}}=\frac{d
\mathbf{r}}{dt} $, \textcolor{Blue}{speed}$=\left| \mathbf{v} \right| $,
\textcolor{Blue}{acceleration} $\mathbf{a}=\frac{d
\mathbf{v}}{dt}=\frac{d^{2}\mathbf{r}}{dt^{2}} $, and the
\textcolor{Blue}{direction of motion} at time $t$ is
$\frac{\mathbf{v}}{\left| \mathbf{v} \right| }$

\paragraph{Integration}
We define the \textbf{indefinite integral} $\int_{{}}^{{}} {\mathbf{r}} \: d{t}
{}$ the set of all anti-derivatives of $\mathbf{r}$.
\[\int_{{}}^{{}} {\mathbf{r}} \: d{t} {}=\left<\int_{{}}^{{}} {f(t)} \: d{t} {}  ,\int_{{}}^{{}} {g(t)} \: d{t} {},\int_{{}}^{{}} {h(t)} \: d{t} {}\right>.\] 
The \textbf{definite integral} is defined, given integrable components over the
interval $\left[ a,b \right]$: $\int_{{a}}^{{b}} {\mathbf{r}} \: d{t} {}=$
\[\left( \int_{{a}}^{{b}} {f(t)}\: d{t} {} \right)\hat{\imath}+\left( \int_{{a}}^{{b}} {g(t)} \: d{t} {} \right)\hat{\jmath  }+ \left( \int_{{a}}^{{b}} {h(t)} \: d{t} {} \right) \hat{k}\]

\textcolor{Blue}{\paragraph{Application: Projectile Motion}}
We have the following results from application of vector calculus on
\textcolor{ForestGreen}{projectile motion}, thrown at an angle $\alpha$ above
the $x$-axis with initial velocity $\mathbf{v}_0$, and $v_0=\left| \mathbf{v}_0
\right| $:
\begin{itemize}
  \item Horizontal displacement $x=tv_0\cos{\alpha}$
  \item Vertical acceleration $y'' = -g$ 
  \item Vertical velocity $\frac{dy}{dt} = -gt + C_1=v_0\sin{\alpha}-gt$
  \item Vertical displacement $y=tv_0\sin{\alpha}-\frac{1}{2}gt^{2}+C_2$
  \item Displacement $\mathbf{r}=\left<rv_0\cos{\alpha},tv_0\sin{\alpha}-\frac{1}{2}gt^{2} \right> $
  \item Velocity $\mathbf{v}=\left<v_0\cos{\alpha},v_0\sin{\alpha}-gt\right> $
  \item Maximum height at $y'=0$, $tmax=\frac{v_0\sin{\alpha}}{g}$
  \item $y=0$ at $t=0$ or $t=\frac{2v_0\sin{\alpha}}{g}$
\end{itemize}

\paragraph{Arc Lengths}
To motivate the derivation of arc length, we consider a short straight segment
of the curve 
\begin{align*}
  \Delta L\approx&\sqrt{{\left( \Delta x \right) }^{2}+{\left(
\Delta y \right) }^{2}+{\left( \Delta z \right) }^{2}}\\
        =&\sqrt{{\left( \frac{dx}{dt} \right) }^{2}+{\left( \frac{dy}{dt} \right) }^{2}+{\left( \frac{dz}{dt} \right) }^{2}}\Delta t
\end{align*}
Taking the sum of the small arc segments, we obtain the length of the smooth
curve
$\mathbf{r}\left( t \right) =\left< x\left( t \right) ,y(t),z(t) \right>,\quad
a\le t\le b$ as $t$ increases from $t=a$ to $t=b$ is:
\[L=\int_{{a}}^{{b}} {\sqrt{{\left( \frac{dx}{dt} \right) }^{2} +{\left( \frac{dy}{dt} \right) }^{2}+{\left( \frac{dz}{dt} \right) }^{2}}} \: d{t} {}.\]
We call $L$ the \textcolor{Bittersweet}{arc length} of the curve.\\
Considering a function $s(t)=\int_{{t_0}}^{{t}} {\left| \mathbf{v}\left( t
\right)  \right| } \: d{t} {}$ the arc length from $t_0$ to $t$, we have the \textcolor{Bittersweet}{unit
tangent vector} of the smooth curve
\[\mathbf{T}\left( t \right)=\frac{\mathbf{v}(t)}{\left| \mathbf{v}(t) \right| }=\frac{d\mathbf{r}(s)}{ds}=\frac{d\mathbf{r}(t)/dt}{ds/dt}.\] 

\paragraph{Curvature}
As a measure of the amount of ``curving'', we define the
\textcolor{Bittersweet}{curvature} of a smooth curve:
\[\kappa=\left| \frac{d \mathbf{T}\left( s \right) }{ds} \right|=\frac{1}{\left| \mathbf{v} \right|} \left| \frac{d \mathbf{T}}{dt} \right|.\]
That is, the rate of change of direction of the unit tangent vector.
We may also interpret curvature as the \textcolor{Blue}{length of the
acceleration vector} while travelling on the curve at $\left| \mathbf{v}
\right| =1$.
